\section{超越次数}\label{sec:transcendence-degree}

上一节熟悉了超越基. 域扩张的超越基在许多方面与向量空间的基类似. 重要的是, 向量空间的基提供了度量向量空间大小的方法: 所有基的大小都一样, 向量空间的维数就是任一基的大小. 本节为超越基发展相应的概念. 具体来说, 我们证明每个有限生成域扩张都有有限超越基, 且所有这样的基大小相同. 于是可以定义域扩张 $K\subseteq L$ 的超越次数, 用代数的方式 (而非仅仅用线性代数的方式) 度量 $L$ 相对于 $K$ 的大小.

先证明有限超越基的存在性.

\begin{Proposition}{有限超越基的存在性}\leavevmode\label{prop:finite-transcendence-basis-existence}
设 $K\subseteq L$ 为域, 且 $a_{1},\ldots,a_{n}\in L$. 若 $L$ 在 $K(a_{1},\ldots,a_{n})$ 上代数, 则 $\{a_{1},\ldots,a_{n}\}$ 包含 $L$ 在 $K$ 上的一个超越基. 特别地, 所有有限生成域扩张都有有限超越基.
\end{Proposition}
\begin{Proof}
对 $n$ 归纳. 若 $n=0$, 则 $L$ 在 $K$ 上代数, 且 $\varnothing$ 是一个超越基 (Example~\ref{ex:empty-transcendence-basis}), 归纳起点成立.

假设结论对少于 $n$ 个元素成立. 取 $a_{1},\ldots,a_{n}\in L$, 使得 $L$ 在 $K(a_{1},\ldots,a_{n})$ 上代数. 若 $a_{1},\ldots,a_{n}$ 代数独立, 则它们组成 $L$ 在 $K$ 上的一个超越基, 结论得证. 否则存在非零多项式 $f\in K[x_{1},\ldots,x_{n}]$, 使得 $f(a_{1},\ldots,a_{n})=0$. 在所有这样的 $f$ 中, 选取一个使 $f$ 关于尽可能少的变量有正次数. 因为 $f$ 至少关于一个 $x_{i}$ 有正次数, 不妨设它关于 $x_{n}$ 有正次数. 若写作 $f=\sum g_{i}x_{n}^{i}$, 其中 $g_{i}\in K[x_{1},\ldots,x_{n-1}]$, 则至少一个 $g_{i}$ 非零. 而且, 任何非零的 $g_{i}$ 在 $(a_{1},\ldots,a_{n-1})$ 处都不能取值为零, 否则就违背了选取 $f$ 时的极小性. 因此, 一元多项式
\[
  f(a_{1},\ldots,a_{n-1},x_{n})\in K[a_{1},\ldots,a_{n-1}][x_{n}]
\]
非零, 且在 $a_{n}$ 处取值为零. 因而 $a_{n}$ 在 $K[a_{1},\ldots,a_{n-1}]$ 上代数, 从而在 $K(a_{1},\ldots,a_{n-1})$ 上代数. 由 Proposition~\ref{prop:extending-fields-by-algebraic-elements},
\[
  K(a_{1},\ldots,a_{n-1})\subseteq K(a_{1},\ldots,a_{n})
\]
是代数域扩张. 现在有一串代数扩张
\[
  K(a_{1},\ldots,a_{n-1})\subseteq K(a_{1},\ldots,a_{n})\subseteq L,
\]
而 Corollary~\ref{cor:algebraic-field-extensions-transitive} 蕴含 $K(a_{1},\ldots,a_{n-1})\subseteq L$ 是代数扩张. 由归纳假设, $\{a_{1},\ldots,a_{n-1}\}$ 包含 $L$ 在 $K$ 上的一个超越基. 所以 $\{a_{1},\ldots,a_{n}\}$ 也包含 $L$ 在 $K$ 上的一个超越基, 归纳论证完成.

为说明引理中的最后一个断言, 假设 $L$ 在 $K$ 上有限生成. 那么对某些 $a_{1},\ldots,a_{n}\in L$, $L$ 等于 $K(a_{1},\ldots,a_{n})$, 因而在其上代数. 所以 $\{a_{1},\ldots,a_{n}\}$ 包含一个有限超越基.
\end{Proof}

\BookPairedNoteLeft
  {接下来的任务是证明每个超越基都有相同的大小. 下面的引理以其证明中使用的策略命名, 是我们所需的关键结论.}
  {再多做一点工作并利用选择公理, 可以证明所有域扩张都存在 (可能无限的) 超越基.}

\begin{Lemma}{交换引理}\leavevmode\label{lem:transcendence-basis-exchange}
设 $K\subseteq L$ 为域, $a_{1},\ldots,a_{n}\in L$ 是 $L$ 中互不相同的元素, 且 $L$ 在 $K(a_{1},\ldots,a_{n})$ 上代数. 若集合 $S\subseteq L$ 在 $K$ 上代数独立, 则 $S$ 有限且 $|S|\leq n$.
\end{Lemma}
\begin{Proof}
设 $a_{1},\ldots,a_{n}\in L$ 为 $L$ 中互不相同的元素, $L$ 在 $K(a_{1},\ldots,a_{n})$ 上代数, 且 $S\subseteq L$ 在 $K$ 上代数独立. 反设 $|S|>n$ (包括 $S$ 无限的可能). 为导出矛盾, 依次把 $\{a_{1},\ldots,a_{n}\}$ 中的元素\textquotedbl{}交换\textquotedbl{}成 $S$ 中的元素. 更具体地说, 完成第 $k$ 步时, 我们已选出了互不相同的 $b_{1},\ldots,b_{k}\in S$, 使得在必要时重新排列 $a_{1},\ldots,a_{n}$ 后, $L$ 在 $K(b_{1},\ldots,b_{k},a_{k+1},\ldots,a_{n})$ 上代数. 如下所示, 矛盾在完成第 $n$ 步后出现. 现在精确描述这个递归交换过程的第 $k$ 步.

设对某个 $k\in\{0,\ldots,n\}$, 已找到互不相同的元素 $b_{1},\ldots,b_{k-1}\in S$, 使得在必要时重新排列 $a_{1},\ldots,a_{n}$ 后, $L$ 在 $K(b_{1},\ldots,b_{k-1},a_{k},\ldots,a_{n})$ 上代数. 由于 $|S|>n\geq k$, 可以再选一个元素 $b_{k}\in S\setminus\{b_{1},\ldots,b_{k-1}\}$. 由于 $L$ 在域 $K(b_{1},\ldots,b_{k-1},a_{k},\ldots,a_{n})$ 上代数, 它也在代数 $K[b_{1},\ldots,b_{k-1},a_{k},\ldots,a_{n}]$ 上代数 (Exercise~\ref{ex:algebraic-over-generated-field-iff-algebraic-over-algebra}). 由于 $b_{k}\in L$, 这意味着存在非零多项式 $f\in K[x_{1},\ldots,x_{k},y_{k},\ldots,y_{n}]$, 使得
\[
  f(b_{1},\ldots,b_{k},a_{k},\ldots,a_{n})=0.
\]
在所有这样的 $f$ 中, 选取一个使 $f$ 关于 $y_{k},\ldots,y_{n}$ 中尽可能少的变量有正次数. 由于 $b_{1},\ldots,b_{k}$ 是代数独立集 $S$ 中互不相同的元素, $f$ 必须关于 $y_{k},\ldots,y_{n}$ 中至少一个变量有正次数. 必要时重新编号, 不妨设 $f$ 关于 $y_{k}$ 有正次数. 把 $f$ 看作系数在 $K[x_{1},\ldots,x_{k},y_{k+1},\ldots,y_{n}]$ 中的 $y_{k}$ 的一元多项式, 注意它的任何非零系数在 $(b_{1},\ldots,b_{k},a_{k+1},\ldots,a_{n})$ 处都不能取值为零, 否则就违背了选取 $f$ 时的极小性. 因而
\[
  \begin{aligned}
    &f(b_{1},\ldots,b_{k},y_{k},a_{k+1},\ldots,a_{n})\\
    &\qquad\in K[b_{1},\ldots,b_{k},a_{k+1},\ldots,a_{n}][y_{k}]
  \end{aligned}
\]
是一个非零多项式, 且在 $a_{k}$ 处取值为零. 所以 $a_{k}$ 在 $K[b_{1},\ldots,b_{k},a_{k+1},\ldots,a_{n}]$ 上代数, 也在 $K(b_{1},\ldots,b_{k},a_{k+1},\ldots,a_{n})$ 上代数. 因此下列链中的每个域都在它前一个域上代数:
\[
  \begin{gathered}
    K(b_{1},\ldots,b_{k},a_{k+1},\ldots,a_{n})\\
    \subseteq K(b_{1},\ldots,b_{k},a_{k},\ldots,a_{n})\subseteq L.
  \end{gathered}
\]
于是由 Corollary~\ref{cor:algebraic-field-extensions-transitive}, $L$ 在 $K(b_{1},\ldots,b_{k},a_{k+1},\ldots,a_{n})$ 上代数, 交换过程的第 $k$ 步完成.

完成上述递归过程的第 $n$ 步后, 得到不同的元素 $b_{1},\ldots,b_{n}\in S$, 使得 $L$ 在 $K(b_{1},\ldots,b_{n})$ 上代数. 按假设, $|S|>n$, 所以存在 $b\in S\setminus\{b_{1},\ldots,b_{n}\}\subseteq L$. 因此 $b$ 在 $K(b_{1},\ldots,b_{n})$ 上代数, 也在 $K[b_{1},\ldots,b_{n}]$ 上代数. 于是有非平凡的多项式关系 $f(b_{1},\ldots,b_{n},b)=0$, 与 $S$ 的代数独立性矛盾. 这一矛盾表明 $S$ 有限且 $|S|\leq n$.
\end{Proof}

由交换引理可推出下列重要结论.

\begin{Proposition}{所有超越基的大小相同}\leavevmode\label{prop:transcendence-bases-equal-size}
若 $L$ 是 $K$ 的有限生成域扩张, 则 $L$ 在 $K$ 上的每个超越基都有限, 且大小相同.
\end{Proposition}
\begin{Proof}
设 $L$ 是 $K$ 的有限生成域扩张. 由 Proposition~\ref{prop:finite-transcendence-basis-existence}, $L$ 包含有限超越基, 取 $\{a_{1},\ldots,a_{n}\}$ 为一个大小为 $n$ 的超越基. 设 $S$ 为任意另一个超越基. 由于 $L$ 在 $K(a_{1},\ldots,a_{n})$ 上代数, 且 $S$ 在 $K$ 上代数独立, 交换引理蕴含 $S$ 有限且 $|S|\leq n$. 写成 $S=\{b_{1},\ldots,b_{m}\}$, 其中 $b_{1},\ldots,b_{m}\in L$ 互不相同且 $m\leq n$. 由于 $L$ 在 $K(b_{1},\ldots,b_{m})$ 上代数, 且 $a_{1},\ldots,a_{n}$ 在 $K$ 上代数独立, 交换引理蕴含 $m\geq n$. 因此 $|S|=m=n$.
\end{Proof}

现在得到下面的核心定义. 它是本章余下部分定义和研究仿射簇维数的关键.

\begin{Definition}{超越次数}\leavevmode\label{def:transcendence-degree}
设 $L$ 是 $K$ 的有限生成域扩张. $L$ 在 $K$ 上的\emph{超越次数}, 记作 $\operatorname{trdeg}_{K}(L)$, 是 $L$ 在 $K$ 上的任一超越基的大小.
\end{Definition}

再来看上一节的几个例子.

\begin{Example}{代数扩张的超越次数}\leavevmode\label{ex:transcendence-degree-of-algebraic-extensions}
从 Example~\ref{ex:empty-transcendence-basis} 回忆, $K\subseteq L$ 是代数扩张当且仅当 $\varnothing$ 是一个超越基. 换言之, $K\subseteq L$ 是代数扩张当且仅当 $\operatorname{trdeg}_{K}(L)=0$. 从这个意义上说, 代数域扩张是最小的一类域扩张.
\end{Example}

\begin{Example}{超越次数与单位球面}\leavevmode\label{ex:transcendence-degree-of-unit-sphere}
令 $K=\mathbb{C}$, 并令 $X=\mathcal{V}(x^{2}+y^{2}+z^{2}-1)\subseteq\mathbb{A}^{3}$. 如 Example~\ref{ex:transcendence-basis-of-unit-sphere} 所见, 二元集合 $\{[x],[y]\}\subseteq K(X)$ 是 $K$ 上的一个超越基. 因此 $\operatorname{trdeg}_{K}K(X)=2$.
\end{Example}

\begin{Example}{$\mathbb{Q}(\pi,e)$ 的超越次数}\leavevmode\label{ex:transcendence-degree-pi-e}
换一种说法表述 Example~\ref{ex:pi-and-e-transcendence-basis} 的未解问题: 在写作本书时, 仍不知道 $\operatorname{trdeg}_{\mathbb{Q}}\mathbb{Q}(\pi,e)$ 是 $1$ 还是 $2$.
\end{Example}

本节最后列出几个可用于估计超越次数的性质. 注意 Proposition~\ref{prop:finite-transcendence-basis-existence} 提供了\textquotedbl{}从上方\textquotedbl{}求超越基和超越次数的方法: 若 $L$ 在 $K(S)$ 上代数, 则可以把 $S\subseteq L$ 缩减成 $L$ 的一个超越基, 类似于将向量空间的生成集缩减成一组基. 下一个结论提供了\textquotedbl{}从下方\textquotedbl{}求超越基和超越次数的方法: 若 $S\subseteq L$ 是代数独立集, 则可以把它扩充成 $L$ 在 $K$ 上的一个超越基.

\begin{Proposition}{独立集包含在基中}\leavevmode\label{prop:independent-sets-extend-to-transcendence-bases}
设 $K\subseteq L$ 为有限生成域扩张. 若子集 $S\subseteq L$ 在 $K$ 上代数独立, 则 $S$ 有限, 且包含在 $L$ 在 $K$ 上的一个超越基中.
\end{Proposition}
\begin{Proof}
见 Exercise~\ref{ex:prove-independent-sets-extend-to-transcendence-bases}.
\end{Proof}

由 Proposition~\ref{prop:finite-transcendence-basis-existence} 和~\ref{prop:independent-sets-extend-to-transcendence-bases}, 得到以下结论, 它提供了确定超越次数上下界的方法.

\begin{Corollary}{超越次数的界}\leavevmode\label{cor:bounds-on-transcendence-degree}
设 $K\subseteq L$ 为有限生成域扩张, $S\subseteq L$ 为子集.
\begin{enumerate}
  \item\label{part:transcendence-degree-upper-bound} 若 $L$ 在 $K(S)$ 上代数, 则 $\operatorname{trdeg}_{K}(L)\leq|S|$.
  \item\label{part:transcendence-degree-lower-bound} 若 $S$ 在 $K$ 上代数独立, 则 $\operatorname{trdeg}_{K}(L)\geq|S|$.
\end{enumerate}
\end{Corollary}
\begin{Proof}
超越次数有限, 所以若 $S$ 无限, 第一个断言为真, 因而可假设 $S$ 有限. 若 $L$ 在 $K(S)$ 上代数, 则 Proposition~\ref{prop:finite-transcendence-basis-existence} 蕴含 $S$ 包含一个超越基. 由超越次数的定义得 $\operatorname{trdeg}_{K}(L)\leq|S|$. 类似地, 若 $S$ 在 $K$ 上代数独立, 则 Proposition~\ref{prop:independent-sets-extend-to-transcendence-bases} 蕴含 $S$ 有限且包含在一个超越基中, 从而 $\operatorname{trdeg}_{K}(L)\geq|S|$.
\end{Proof}

\BookExerciseHeading
\begin{BookExercises}
  \item\label{ex:prove-independent-sets-extend-to-transcendence-bases} 证明 Proposition~\ref{prop:independent-sets-extend-to-transcendence-bases}.
  \item\label{ex:transcendence-degree-monotone-in-field} 设 $L\subseteq M$ 是在 $K$ 上有限生成的域. 证明 $\operatorname{trdeg}_{K}(L)\leq\operatorname{trdeg}_{K}(M)$.
  \item\label{ex:transcendence-degree-of-pi-i-square-root-two} 求 $\mathbb{Q}(\pi,i,\sqrt{2})$ 在 $\mathbb{Q}$ 上的超越次数.
  \item\label{ex:characterize-basis-of-known-transcendence-degree} 设 $K\subseteq L$ 是有限生成域扩张, $\operatorname{trdeg}_{K}(L)=d$, 且 $a_{1},\ldots,a_{d}\in L$ 互不相同.
  \begin{enumerate}
    \item 证明 $\{a_{1},\ldots,a_{d}\}\subseteq L$ 是 $L$ 在 $K$ 上的一个超越基, 当且仅当 $a_{1},\ldots,a_{d}$ 在 $K$ 上代数独立.
    \item 证明 $\{a_{1},\ldots,a_{d}\}\subseteq L$ 是 $L$ 在 $K$ 上的一个超越基, 当且仅当 $L$ 在 $K(a_{1},\ldots,a_{d})$ 上代数.
  \end{enumerate}
  \item\label{ex:extremal-characterizations-of-transcendence-degree} 设 $K\subseteq L$ 为有限生成域扩张.
  \begin{enumerate}
    \item 证明 $\operatorname{trdeg}_{K}(L)$ 是 $L$ 中代数独立元素的最大个数.
    \item 证明 $\operatorname{trdeg}_{K}(L)$ 是满足以下条件的 $L$ 中元素的最小个数: 它们生成的域使 $L$ 在该域上代数.
  \end{enumerate}
  \item\label{ex:linear-algebra-analogues-of-transcendence-degree} 用向量空间的维数, 张成空间和线性独立等概念, 陈述本节所有结论以及 Exercise~\ref{ex:characterize-basis-of-known-transcendence-degree} 和~\ref{ex:extremal-characterizations-of-transcendence-degree} 的线性代数对应结论.
\end{BookExercises}
